% Figure 4: the Pareto view. Solved tasks against dollars spent - the frontier
% that decides which configuration a practitioner should actually run.
% Coordinates from figures/data/pareto.dat.

\begin{figure}[t]
\centering
\begin{tikzpicture}[font=\small]
\begin{axis}[
  width=0.92\linewidth, height=5.2cm,
  xmin=10.5, xmax=19.5, ymin=75, ymax=180,
  xlabel={стоимость серии (USD, меньше — лучше)},
  ylabel={решённых задач (больше — лучше)},
  grid=major, grid style={black!12},
  nodes near coords={\scriptsize\pgfplotspointmeta},
  point meta=explicit symbolic,
  every node near coord/.append style={anchor=south, yshift=3pt},
]
\addplot[only marks, mark=*, mark size=3pt, draw=black!70, fill=blue!45]
  table[x=cost, y=solved, meta=label] {figures/data/pareto.dat};

\draw[->, thick, black!45] (axis cs:18.6,96) -- (axis cs:17.4,96);
\node[font=\scriptsize, text=black!55, anchor=west] at (axis cs:18.4,90)
  {лучше};
\draw[->, thick, black!45] (axis cs:18.9,98) -- (axis cs:18.9,110);
\end{axis}
\end{tikzpicture}
\caption{Граница, которая имеет значение на практике: решённые задачи против
потраченных денег, вся серия. Опубликованная конфигурация ($k{=}0$)
доминируема: она тратит меньше контрольной группы, но решает намного меньше
задач, и обе $k{=}2$ и $k{=}6$ превосходят её по обеим осям одновременно.
Читайте этот рисунок только как обобщение на уровне серии. Поскольку
неуспешный запуск останавливается раньше и потому стоит немного, итоги по
серии приукрашивают руки, которые проваливаются чаще, поэтому именно
сравнение при совпадающем исходе из Таблицы~\ref{tab:cost}, а не этот
рисунок, несёт утверждение статьи о стоимости. По этому честному сравнению
$k{=}6$ \emph{не} дешевле контрольной группы (\CostRatioTailSix$\times$,
$p={}$\CostSignPTailSix). Что рисунок действительно показывает — это то, что
ни один из крайних случаев оси хвоста не является правильным ответом.}
\label{fig:pareto}
\end{figure}
